\documentclass[12pt,a4paper]{article}

% ── Packages ─────────────────────────────────────────────────────────────────
\usepackage[margin=2.5cm]{geometry}
\usepackage{amsmath,amssymb,amsthm}
\usepackage{graphicx}
\usepackage{booktabs}
\usepackage{hyperref}
\usepackage{xcolor}
\usepackage{listings}
\usepackage{caption}
\usepackage{float}
\usepackage{microtype}
\usepackage{enumitem}
\usepackage{mdframed}
\usepackage{tocloft}

% ── Hyperref colours (match existing report) ─────────────────────────────────
\hypersetup{
  colorlinks=true, linkcolor=blue!60!black,
  citecolor=green!50!black, urlcolor=blue!70!black,
}

% ── Python listings style (match existing report) ────────────────────────────
\lstset{
  language=Python,
  basicstyle=\ttfamily\small,
  keywordstyle=\color{blue!70!black}\bfseries,
  commentstyle=\color{gray},
  stringstyle=\color{green!50!black},
  breaklines=true,
  frame=single,
  rulecolor=\color{gray!40},
  backgroundcolor=\color{gray!5},
  numbers=left, numberstyle=\tiny\color{gray},
  numbersep=6pt,
  literate={←}{{$\leftarrow$}}1 {→}{{$\rightarrow$}}1
           {∂}{{$\partial$}}1 {≈}{{$\approx$}}1
           {‖}{{$\|$}}1,
}

% Plain listings (pseudocode, no syntax highlighting)
\lstdefinestyle{plain}{
  language={},
  basicstyle=\ttfamily\small,
  breaklines=true,
  frame=single,
  rulecolor=\color{gray!40},
  backgroundcolor=\color{gray!5},
  numbers=none,
}

% ── Custom commands ───────────────────────────────────────────────────────────
\newcommand{\norm}[1]{\left\lVert#1\right\rVert}
\newcommand{\Lap}{\Delta}
\newcommand{\lapinv}{\Delta^{-1}}
\newcommand{\hminusone}{H^{-1}}
\newcommand{\code}[1]{\texttt{#1}}
% Code identifiers that appear in display math
\newcommand{\LAPINV}{\text{\texttt{LAP\_INV}}}
\newcommand{\LAMBDAINV}{\text{\texttt{LAMBDA\_INV}}}
\newcommand{\DELX}{\text{\texttt{DEL\_X}}}
\newcommand{\DELY}{\text{\texttt{DEL\_Y}}}

% Highlight box for key insights
\newmdenv[
  backgroundcolor=blue!4,
  linecolor=blue!40,
  linewidth=1pt,
  innertopmargin=6pt,
  innerbottommargin=6pt,
  innerleftmargin=10pt,
  innerrightmargin=10pt,
  skipabove=8pt,
  skipbelow=8pt,
]{insightbox}

\theoremstyle{definition}
\newtheorem{remark}{Remark}

% ── Title ─────────────────────────────────────────────────────────────────────
\title{%
  \textbf{A Beginner's Guide to the Discrete Fourier Transform\\
  and Spectral Methods}\\[8pt]
  \large For Readers of \textit{Optimal Mixing of Passive Scalars}\\[4pt]
  \normalsize Python Implementation and Numerical Results
}
\author{Adebanji Adelowo}
\date{\today}

% ─────────────────────────────────────────────────────────────────────────────
\begin{document}
\maketitle
\tableofcontents
\bigskip\hrule\bigskip

% ── Preface ───────────────────────────────────────────────────────────────────
\section*{Why This Guide Exists}
\addcontentsline{toc}{section}{Why This Guide Exists}

The Numerical Methods section (Section~3) of the optimal mixing report uses
concepts from Fourier analysis and spectral methods — tools from signal
processing and applied mathematics that can feel very abstract at first.
This guide explains every concept from the ground up.

\begin{insightbox}
\textbf{If you have studied calculus and basic linear algebra, you have
everything you need to follow this guide.}
\end{insightbox}

% ─────────────────────────────────────────────────────────────────────────────
\section{Part 1: The Big Picture — What Is a Spectral Method?}

\subsection{The Problem We Want to Solve}

The mixing simulation involves solving a partial differential equation (PDE):
\begin{equation}
  \frac{\partial \theta}{\partial t} + u \cdot \nabla \theta = 0.
\end{equation}
This describes how a scalar field $\theta(x,y,t)$ — think of it as the
concentration of ink in water — changes over time under a velocity field $u$.
The term $\nabla\theta$ means ``gradient of $\theta$'': a vector pointing in
the direction of steepest increase.

To solve this on a computer we must \emph{discretise} it: approximate it on a
finite grid of points.  There are many ways to do this.  The
\textbf{spectral method} is one of the most accurate.

\subsection{Two Ways to Represent a Function}

Imagine a wiggly function $f(x)$ defined on $[0,1]$:

\begin{description}
  \item[Way 1 — Physical space.] Store the values
    $f(x_0), f(x_1), \ldots, f(x_{N-1})$ at grid points.  Think of it as a
    bar chart of heights.
  \item[Way 2 — Frequency space (Fourier space).] Represent $f$ as a sum of
    sine and cosine waves of different frequencies.  Instead of storing
    heights, store \emph{how much of each frequency is present}.
\end{description}

These are just two descriptions of the same function — like describing a song
by its audio waveform versus its sheet music.  You can switch between them
freely.

\begin{insightbox}
\textbf{Key insight:} Some operations that are hard in physical space (like
computing derivatives) become trivially simple in frequency space — just
multiply by a number.
\end{insightbox}

% ─────────────────────────────────────────────────────────────────────────────
\section{Part 2: Waves and Frequencies}

\subsection{What Is a Frequency?}

A pure wave of frequency $k$ on $[0,1]$ looks like:
\[
  f_k(x) = \sin(2\pi k x) \qquad \text{or} \qquad f_k(x) = \cos(2\pi k x).
\]
The number $k$ counts how many complete oscillations fit inside $[0,1]$:
\begin{itemize}
  \item $k=1$: one full oscillation — a single gentle curve.
  \item $k=2$: two full oscillations — a slightly wigglier curve.
  \item $k=10$: ten oscillations — a rapidly oscillating wave.
\end{itemize}
Higher $k$ = higher frequency = finer (smaller-scale) features.

\subsection{Any Periodic Function Is a Sum of Waves}

This is the central claim of Fourier analysis:
\textbf{any periodic function can be built by adding up sine and cosine waves}.
This sum is the \textbf{Fourier series}:
\begin{equation}
  f(x) = a_0 + \sum_{k=1}^{\infty}\bigl[a_k\cos(2\pi kx) + b_k\sin(2\pi kx)\bigr].
\end{equation}
The coefficients $a_k$ and $b_k$ tell you how much of each frequency is present.

\paragraph{A concrete example.}
A square wave (alternating $+1$ and $-1$) can be written as:
\[
  f(x) = \frac{4}{\pi}\!\left[\sin(2\pi x)
    + \frac{1}{3}\sin(6\pi x)
    + \frac{1}{5}\sin(10\pi x) + \cdots\right].
\]
Adding the first few terms already gives a recognisable approximation; the
infinite series is exact.

\begin{figure}[H]
  \centering
  \includegraphics[width=\linewidth]{figures/fig_fourier_series.pdf}
  \caption{Left: the first three sine-wave components of the square wave's Fourier
    series.  Right: adding more and more terms ($1$, $3$, $9$) progressively
    improves the approximation.  The Gibbs overshoot near the jumps shrinks as
    more terms are included but never entirely disappears.}
  \label{fig:fourier_series}
\end{figure}

\subsection{Complex Exponentials — A Compact Notation}

Instead of writing $\sin$ and $\cos$ separately, mathematicians use
\textbf{complex exponentials} via Euler's formula:
\[
  e^{2\pi ikx} = \cos(2\pi kx) + i\sin(2\pi kx).
\]
The Fourier series then takes the symmetric form:
\begin{equation}
  f(x) = \sum_{k=-\infty}^{\infty} \hat{f}(k)\,e^{2\pi ikx}.
\end{equation}
The complex number $\hat{f}(k)$ is the \textbf{Fourier coefficient} at
wavenumber $k$.  Its magnitude $|\hat{f}(k)|$ is the amplitude of that
frequency, and its argument $\arg(\hat{f}(k))$ is the phase shift.

% ─────────────────────────────────────────────────────────────────────────────
\section{Part 3: The Discrete Fourier Transform (DFT)}

\subsection{From Continuous to Discrete}

The Fourier series works for continuous functions.  Computers work with finite
lists of numbers.  We need a discrete version.

Suppose we have $N$ equally spaced sample points:
\[
  x_j = \frac{j}{N}, \quad j = 0, 1, 2, \ldots, N-1.
\]
The \textbf{Discrete Fourier Transform (DFT)} takes the $N$ values
$f[0], f[1], \ldots, f[N-1]$ and produces $N$ Fourier coefficients:
\begin{equation}
  \hat{f}[k] = \sum_{j=0}^{N-1} f[j]\,e^{-2\pi ijk/N},
  \quad k = 0, 1, \ldots, N-1.
\end{equation}
The \textbf{Inverse DFT (IDFT)} goes back:
\begin{equation}
  f[j] = \frac{1}{N}\sum_{k=0}^{N-1}\hat{f}[k]\,e^{2\pi ijk/N}.
\end{equation}

\begin{insightbox}
\textbf{No information is lost.}  Going forward and back is exact.  The DFT
and IDFT are just two different representations of the same data — like
converting between Celsius and Fahrenheit.
\end{insightbox}

\begin{figure}[H]
  \centering
  \includegraphics[width=\linewidth]{figures/fig_dft_spectrum.pdf}
  \caption{Left: a signal composed of three sine waves at $k=3$, $7$, and $13$.
    Right: the magnitude of its DFT reveals exactly those three frequencies as
    sharp peaks — the DFT has ``heard'' each component separately, even though
    they are all jumbled together in physical space.}
  \label{fig:dft_spectrum}
\end{figure}

\subsection{What Do the Fourier Coefficients Mean?}

$\hat{f}[k]$ is a complex number whose:
\begin{itemize}
  \item \textbf{Magnitude} $|\hat{f}[k]|$ tells you how much of frequency $k$
    is present in $f$.
  \item \textbf{Phase} $\arg(\hat{f}[k])$ tells you how that wave is shifted
    (its offset).
\end{itemize}
If $f$ is real-valued (as physical quantities always are), the
negative-frequency coefficients are redundant:
$\hat{f}[-k] = \overline{\hat{f}[k]}$ (complex conjugate).

\subsection{The FFT — The Practical Tool}

Computing the DFT directly requires $O(N^2)$ operations — for $N=1000$, that
is $10^6$ multiplications.  The \textbf{Fast Fourier Transform (FFT)} algorithm
computes the \emph{exact same result} in $O(N\log N)$ operations — for
$N=1000$, that is only $\approx 10{,}000$.

\medskip
\noindent\emph{FFT vs.\ DFT: the algorithm is different, the result is
identical.}

\begin{lstlisting}
f_hat  = np.fft.fft(f)        # 1D FFT
f_back = np.fft.ifft(f_hat)   # 1D inverse FFT

F_hat  = np.fft.fft2(F)       # 2D FFT (applied to each axis)
F_back = np.fft.ifft2(F_hat)  # 2D inverse FFT
\end{lstlisting}

\subsection{The 2D DFT}

For a 2D function $f(x,y)$ defined on an $N\times N$ grid, the DFT is applied
in both directions independently:
\begin{equation}
  \hat{f}[k_y,k_x]
  = \sum_{i=0}^{N-1}\sum_{j=0}^{N-1} f[i,j]\,e^{-2\pi i(ik_y+jk_x)/N}.
\end{equation}

\paragraph{Index convention used in this project (and NumPy).}
\begin{itemize}
  \item Row index $i$ corresponds to the \textbf{$y$-direction}.
  \item Column index $j$ corresponds to the \textbf{$x$-direction}.
  \item So \code{theta\_hat[ky, kx]} is the coefficient for wavenumber
    $(k_x, k_y)$.
\end{itemize}
This is why \code{DEL\_X} has shape \code{(1,~N)} (acts on columns) and
\code{DEL\_Y} has shape \code{(N,~1)} (acts on rows).

% ─────────────────────────────────────────────────────────────────────────────
\section{Part 4: Wavenumbers — What Are They?}

\subsection{The Wavenumber Array}

After computing \code{fft(f)} on an array of length $N$, the output has $N$
entries corresponding to $N$ different frequencies stored in a specific order.

For $N=8$, the array indices and their corresponding effective wavenumbers are:

\begin{center}
\begin{tabular}{lcccccccc}
  \toprule
  Index      & 0 & 1 & 2 & 3 & 4   & 5  & 6  & 7  \\
  \midrule
  Wavenumber & 0 & 1 & 2 & 3 & $\pm4$ & $-3$ & $-2$ & $-1$ \\
  \bottomrule
\end{tabular}
\end{center}

\begin{itemize}
  \item \textbf{Index 0}: the constant (DC) component, $k=0$ — the average
    value of $f$.
  \item \textbf{Indices 1 to $N/2-1$}: positive frequencies.
  \item \textbf{Index $N/2$}: the Nyquist frequency $k=N/2$ (special case,
    see below).
  \item \textbf{Indices $N/2+1$ to $N-1$}: negative frequencies.
\end{itemize}

In the code:
\begin{lstlisting}
k     = np.arange(N)
k_eff = k - N * (k > N // 2)   # convert index to effective wavenumber
\end{lstlisting}

This is called ``centring'' the wavenumbers.  To visualise a 2D spectrum with
zero in the middle, use \code{np.fft.fftshift}.

\begin{figure}[H]
  \centering
  \includegraphics[width=\linewidth]{figures/fig_wavenumber_order.pdf}
  \caption{Wavenumber storage for $N=16$.  Left: the raw output of
    \texttt{np.fft.fft}.  Indices $0$--$7$ hold positive frequencies; index
    $8$ is the Nyquist; indices $9$--$15$ hold negative frequencies (wrapped
    around).  Right: after \texttt{np.fft.fftshift} the zero-frequency bin
    sits in the middle, which is more intuitive for visualisation.}
  \label{fig:wavenumber_order}
\end{figure}

\subsection{Why Negative Frequencies?}

Negative frequencies arise naturally from the complex exponential formulation.
For a real-valued function $f$, the $k$ and $-k$ modes always come in
complex-conjugate pairs.  Together they produce a real-valued signal:
\[
  \hat{f}[k]\,e^{2\pi ikx} + \hat{f}[-k]\,e^{-2\pi ikx}
  = 2|\hat{f}[k]|\cos(2\pi kx + \phi).
\]
So negative frequencies represent the ``mirror image'' of positive frequencies
— they are not physically different, just part of the mathematical convention.

\subsection{The Nyquist Frequency — The Highest Representable Frequency}

The grid has spacing $dx = 1/N$.  The finest feature representable is a wave
that alternates sign at every grid point — one oscillation per $2\,dx$.  That
frequency is:
\[
  k_{\text{Nyquist}} = \frac{N}{2}.
\]
\textbf{You cannot represent frequencies higher than $N/2$ on a grid with $N$
points.}  If a function has fine-scale features at $k > N/2$, the DFT
\emph{cannot see them} — they get wrapped around and appear as lower-frequency
noise.  This is called \textbf{aliasing}.

\begin{insightbox}
\textbf{Real-world analogy.}  A movie camera shoots at 24~frames per second.
A car wheel spinning at more than 12~revolutions per second appears to spin
\emph{backwards} in the footage — this is aliasing.  The camera cannot
represent frequencies faster than 12~Hz.
\end{insightbox}

\begin{figure}[H]
  \centering
  \includegraphics[width=0.85\linewidth]{figures/fig_aliasing.pdf}
  \caption{Aliasing on a grid with $N=16$ points (Nyquist limit $= 8$).
    The wave $k=3$ (solid blue) and the wave $k=13$ (dashed red) produce
    \emph{identical} values at every grid point (filled circles), so the
    DFT cannot distinguish them.  Any frequency above the Nyquist is
    silently folded back into a lower frequency.}
  \label{fig:aliasing}
\end{figure}

\subsection{The Nyquist Mode and Why It Is Zeroed for First Derivatives}

At $k = N/2$, the wave $e^{2\pi i(N/2)x}$ sampled at grid points
$x_j = j/N$ gives:
\[
  e^{2\pi i(N/2)(j/N)} = e^{\pi ij} = (-1)^j.
\]
This alternates $+1,-1,+1,-1,\ldots$ — the same as $e^{-2\pi i(N/2)x}$.
So $k = +N/2$ and $k = -N/2$ are \textbf{identical on the grid} — they are
ambiguous.

\begin{itemize}
  \item For the \textbf{Laplacian} (second derivative), this is fine:
    $(-k)^2 = k^2$, so the sign does not matter.
  \item For \textbf{first derivatives}, the sign matters: $+k$ and $-k$
    give derivatives with opposite signs.  Since the mode is ambiguous, we
    zero it out:
\end{itemize}
\begin{lstlisting}
k_d = k_eff.copy()
k_d[N//2] = 0.0   # zero the Nyquist mode for first derivatives
\end{lstlisting}
\code{DEL\_X} and \code{DEL\_Y} use \code{k\_d}; \code{LAP\_INV} uses
\code{k\_eff}.

% ─────────────────────────────────────────────────────────────────────────────
\section{Part 5: Spectral Differentiation — The Core Magic Trick}

\subsection{Why Differentiation Becomes Multiplication}

This is the single most important idea in spectral methods.  Consider a
function written as a sum of waves:
\[
  f(x) = \sum_k \hat{f}(k)\,e^{2\pi ikx}.
\]
Take the derivative term by term:
\[
  \frac{df}{dx}
  = \sum_k \hat{f}(k)\cdot\frac{d}{dx}\!\left[e^{2\pi ikx}\right]
  = \sum_k \hat{f}(k)\cdot(2\pi ik)\,e^{2\pi ikx}.
\]

\begin{insightbox}
\textbf{In Fourier space, taking a derivative simply means multiplying each
coefficient by $2\pi ik$.}
\end{insightbox}

\paragraph{Recipe for computing $df/dx$ spectrally.}
\begin{enumerate}
  \item Compute $\hat{f} = \mathrm{FFT}(f)$.
  \item Multiply: $\widehat{df/dx}[k] = (2\pi ik)\,\hat{f}[k]$.
  \item Transform back: $df/dx = \mathrm{real}\bigl(\mathrm{IFFT}(\widehat{df/dx})\bigr)$.
\end{enumerate}
\begin{lstlisting}
f_hat  = np.fft.fft(f)
df_hat = (2j * np.pi * k_d) * f_hat     # k_d has Nyquist zeroed
df_dx  = np.real(np.fft.ifft(df_hat))   # discard tiny numerical imaginary part
\end{lstlisting}

\paragraph{Accuracy.}
Spectral differentiation is \textbf{spectrally accurate} — the error decreases
faster than any polynomial in $N$ as $N\to\infty$ (for smooth periodic
functions).  It is far more accurate than finite differences, which are only
$O(dx^p)$ accurate for some fixed order $p$.

\begin{figure}[H]
  \centering
  \includegraphics[width=\linewidth]{figures/fig_spectral_derivative.pdf}
  \caption{Left: spectral differentiation (blue) matches the exact derivative
    perfectly at $N=16$, while second-order finite differences (red) show
    visible error.  Right: the convergence plot shows that spectral errors
    drop to machine precision by $N=16$ for this smooth function, whereas
    finite-difference errors decay only as $O(N^{-2})$.  For rough functions
    the spectral advantage is smaller, but for the smooth fields in the mixing
    simulation it is enormous.}
  \label{fig:spectral_derivative}
\end{figure}

\subsection{Second Derivatives and the Laplacian}

Two derivatives each contribute a factor of $2\pi ik$:
\[
  (2\pi ik)^2 = -4\pi^2 k^2.
\]
So the second-derivative multiplier is $-4\pi^2 k^2$ (real and negative).

In 2D, the Laplacian $\Delta f = \partial_{xx}f + \partial_{yy}f$ becomes:
\begin{equation}
  \widehat{\Delta f}[k_y,k_x]
  = \bigl[-(2\pi k_x)^2-(2\pi k_y)^2\bigr]\,\hat{f}[k_y,k_x]
  = -4\pi^2(k_x^2+k_y^2)\,\hat{f}[k_y,k_x].
\end{equation}

\subsection{The Operators \code{DEL\_X} and \code{DEL\_Y} in the Code}

\begin{lstlisting}
DEL_X = (2j * np.pi * k_d)[np.newaxis, :]   # shape (1, N)
DEL_Y = (2j * np.pi * k_d)[:, np.newaxis]   # shape (N, 1)
\end{lstlisting}

For a 2D field $\theta$:
\begin{itemize}
  \item \code{DEL\_X * theta\_hat} multiplies each \emph{column}
    ($x$-direction) by the corresponding $2\pi ik_x$ — giving
    $\widehat{\partial_x\theta}$.
  \item \code{DEL\_Y * theta\_hat} multiplies each \emph{row}
    ($y$-direction) by $2\pi ik_y$ — giving $\widehat{\partial_y\theta}$.
\end{itemize}

\paragraph{Why shapes $(1,N)$ and $(N,1)$?}
NumPy broadcasting means \code{(1,N)} is automatically repeated along rows to
become \code{(N,N)} when multiplied with an \code{(N,N)} array.  This saves
memory without storing any redundant data.

% ─────────────────────────────────────────────────────────────────────────────
\section{Part 6: The Laplacian and Its Inverse}

\subsection{\code{LAP\_INV} — The Inverse Laplacian}

We established that the Laplacian in Fourier space is multiplication by
$-4\pi^2|k|^2$.  The \textbf{inverse Laplacian} is therefore division by the
same factor:
\begin{equation}
  \LAPINV[k_y,k_x]
  = \begin{cases}
      \dfrac{-1}{(2\pi)^2(k_x^2+k_y^2)}, & (k_x,k_y)\neq(0,0),\\[6pt]
      0, & (k_x,k_y)=(0,0).
    \end{cases}
\end{equation}

\paragraph{To apply $\lapinv$ to a field $\theta$:}
\begin{enumerate}
  \item Compute $\hat\theta = \mathrm{FFT2}(\theta)$.
  \item Multiply: $\widehat{\lapinv\theta} = \LAPINV\cdot\hat\theta$.
  \item Transform back: $\lapinv\theta = \mathrm{real}
    \bigl(\mathrm{IFFT2}(\widehat{\lapinv\theta})\bigr)$.
\end{enumerate}

\paragraph{Why zero the DC mode?}
At $k=(0,0)$ we would be dividing by zero.  The DC mode represents the spatial
mean $\langle\theta\rangle = \int\theta\,dx\,dy$.  This mean is conserved and
decoupled from the rest of the dynamics, so we set it to zero.

\paragraph{Physical interpretation of $\lapinv$.}
The inverse Laplacian \textbf{smooths} a function.  High-frequency
(small-scale) features are suppressed by the large $|k|^2$ in the denominator.
If $\theta$ is a concentrated blob, $\lapinv\theta$ is a broad, smooth
distribution.

\subsection{Why We Need $\lapinv$ for the Velocity}

In the LTD optimal velocity formula:
\[
  v = -\lapinv P\!\left(\theta\,\nabla\lapinv\theta\right),
\]
the $\lapinv$ appears \textbf{twice}:
\begin{enumerate}
  \item The inner $\lapinv\theta$ converts the scalar into a
    stream-function-like field (smoother, large-scale).
  \item The outer $-\lapinv$ converts the nonlinear term into the velocity
    itself.
\end{enumerate}
This ensures the velocity has the right structure: divergence-free with
controlled enstrophy.

% ─────────────────────────────────────────────────────────────────────────────
\section{Part 7: Parseval's Identity — Connecting Norms}

\subsection{The Statement}

For a 1D DFT with $N$ points:
\[
  \sum_{j=0}^{N-1}|f[j]|^2 = \frac{1}{N}\sum_{k=0}^{N-1}|\hat{f}[k]|^2.
\]
In the NumPy convention (non-unitary), $\|\hat{f}\|^2 = N\|f\|^2$, so:
\[
  \|f\|_{L^2} \approx \frac{\|\hat{f}\|_{\ell^2}}{N}.
\]
In 2D (with $N^2$ points):
\begin{equation}
  \|f\|_{L^2} \approx \frac{\|\hat{f}\|_{\ell^2}}{N^2}.
\end{equation}

\subsection{What It Means}

Parseval's identity says: \textbf{the total energy (sum of squares) is the
same in physical space and Fourier space}, up to a normalisation factor.

This is analogous to the Pythagorean theorem: if you rotate a vector, its
length does not change.  The DFT is a rotation in an $N$-dimensional space, so
it preserves lengths.

\subsection{How It Is Used in the Code}

To compute $\|\nabla v\|_{L^2}$ (needed for the enstrophy normalisation), the
code computes in Fourier space:
\begin{lstlisting}
grad_v_norm = (1/N**2) * np.sqrt(
    np.sum(np.abs(DEL_X * v_x_hat)**2) +
    np.sum(np.abs(DEL_Y * v_x_hat)**2) +
    np.sum(np.abs(DEL_X * v_y_hat)**2) +
    np.sum(np.abs(DEL_Y * v_y_hat)**2)
)
\end{lstlisting}
Each \code{DEL\_X * v\_x\_hat} is $\widehat{\partial_x v_x}$ in Fourier space,
and the Parseval factor $1/N^2$ converts the Fourier-space norm to the
physical-space $L^2$ norm.

% ─────────────────────────────────────────────────────────────────────────────
\section{Part 8: The \texorpdfstring{$H^{-1}$}{H\^{}\{-1\}} Mix Norm}

\subsection{Why Not Just Use the $L^2$ Norm?}

Under the passive scalar equation $\partial_t\theta + u\cdot\nabla\theta = 0$
with divergence-free $u$, all $L^p$ norms are \textbf{conserved}:
\[
  \|\theta(t)\|_{L^2} = \|\theta_0\|_{L^2} \quad \text{for all } t.
\]
So the $L^2$ norm cannot tell us anything about mixing — it is constant
regardless of how well the dye has been stirred.  We need a norm that
\emph{decreases} as the scalar gets mixed into finer and finer structures.

\subsection{The $H^{-1}$ Norm}

The \textbf{$\hminusone$ norm} is defined as:
\begin{equation}
  \|\theta\|_{H^{-1}}^2
  = \sum_{k\neq 0}\frac{|\hat\theta(k)|^2}{4\pi^2|k|^2}.
\end{equation}
Compare with the $L^2$ norm: $\|\theta\|_{L^2}^2 = \sum_k|\hat\theta(k)|^2$.

The difference: the $\hminusone$ norm divides each mode's contribution by
$|k|^2$, giving \textbf{much less weight to high-frequency (fine-scale)
content}.

\paragraph{What happens when we stir?}
Initially, the scalar's energy is concentrated at low wavenumbers (smooth,
large blobs).  As the velocity stirs it, energy is transferred to higher and
higher wavenumbers (finer filaments).  The $\hminusone$ norm penalises low-$k$
content heavily — so as energy moves to high $k$, the norm decreases.

\subsection{Intuition for the $H^{-1}$ Norm}

Think of it this way:
\begin{itemize}
  \item A scalar with all its energy at low $k$ (one big blob) has
    \textbf{large} $\hminusone$ norm — it is unmixed.
  \item A scalar with its energy spread uniformly across many $k$ values has
    \textbf{small} $\hminusone$ norm — it is well-mixed.
  \item A perfectly uniform scalar ($\theta = \mathrm{const}$, all energy at
    $k=0$) has zero $\hminusone$ norm — perfectly mixed.
\end{itemize}

\subsection{The $\lambda^{-1}$ Weight}

In code, the mix norm is computed via \code{LAMBDA\_INV}:
\begin{equation}
  \LAMBDAINV[k] = \frac{1}{2\pi|k|} = \sqrt{-\LAPINV[k]},
\end{equation}
\begin{equation}
  \|\theta\|_{H^{-1}}^2
  = \frac{\|\LAMBDAINV\cdot\hat\theta\|_{\ell^2}^2}{N^4}.
\end{equation}
\begin{lstlisting}
LAMBDA_INV  = np.sqrt(np.where(LAP_INV != 0, -LAP_INV, 0.0))
mix_norm_sq = np.sum(np.abs(LAMBDA_INV * theta_hat)**2) / N**4
\end{lstlisting}

\begin{figure}[H]
  \centering
  \includegraphics[width=\linewidth]{figures/fig_hminus1_norm.pdf}
  \caption{The $H^{-1}$ norm as a measure of mixing quality.  Left: an
    unmixed Gaussian blob — all energy concentrated at low wavenumbers,
    $\|\theta\|_{H^{-1}}$ is large.  Centre: a fine-scale chequerboard
    pattern with the \emph{same} $L^2$ norm — energy spread to high
    wavenumbers, $\|\theta\|_{H^{-1}}$ is much smaller.  Right: the
    energy spectrum confirms the shift from low-$k$ to high-$k$ content.
    As the mixing simulation runs, the $H^{-1}$ norm decreases monotonically.}
  \label{fig:hminus1}
\end{figure}

% ─────────────────────────────────────────────────────────────────────────────
\section{Part 9: The Leray Projection — Enforcing Incompressibility}

\subsection{What Is Incompressibility?}

The velocity field $u = (u_x, u_y)$ must satisfy $\nabla\cdot u = 0$:
\[
  \frac{\partial u_x}{\partial x} + \frac{\partial u_y}{\partial y} = 0.
\]
This is the \textbf{incompressibility condition}: fluid is not created or
destroyed anywhere.  In physical terms, any fluid flowing into a region must
flow out equally.

\subsection{The Helmholtz Decomposition}

Any vector field $g$ can be split into two parts:
\[
  g = \underbrace{g_{\text{div-free}}}_{\nabla\cdot g_{\text{div-free}}=0}
    + \underbrace{\nabla\phi}_{\text{pure gradient}},
\]
where $\phi$ is some scalar potential.  The \textbf{Leray projection} extracts
just the divergence-free part:
\[
  Pg = g - \nabla\phi, \qquad \phi = \lapinv(\nabla\cdot g).
\]

\paragraph{Step by step:}
\begin{enumerate}
  \item Compute the divergence: $D = \nabla\cdot g = \partial_x g_x + \partial_y g_y$.
  \item Solve $\Lap\phi = D$, i.e.\ $\phi = \lapinv D$.
  \item Subtract the gradient: $Pg = g - \nabla\phi$.
\end{enumerate}

\subsection{In Fourier Space — It's Just Multiplication}

Everything becomes element-wise multiplication in Fourier space:
\begin{equation}
  \widehat{Pg_x}[k]
  = \hat g_x[k]
  - \DELX[k]\cdot\LAPINV[k]\cdot\bigl(\DELX[k]\cdot\hat g_x[k]
    + \DELY[k]\cdot\hat g_y[k]\bigr).
\end{equation}
In code:
\begin{lstlisting}
div_g_hat = DEL_X * gx_hat + DEL_Y * gy_hat        # divergence of g
Pgx_hat   = gx_hat - DEL_X * LAP_INV * div_g_hat   # project x-component
Pgy_hat   = gy_hat - DEL_Y * LAP_INV * div_g_hat   # project y-component
\end{lstlisting}
No iteration required — just a handful of array multiplications.

% ─────────────────────────────────────────────────────────────────────────────
\section{Part 10: The Pseudo-Spectral Method}

\subsection{The Problem with Products}

A purely spectral method would compute \emph{everything} in Fourier space.
But the term $\theta\cdot\nabla\lapinv\theta$ involves a
\textbf{product of two functions in physical space}.

In Fourier space, a product becomes a \textbf{convolution}:
\[
  \widehat{f\cdot g}[k] = \sum_m \hat f[m]\,\hat g[k-m].
\]
Computing this naively costs $O(N^2)$ per wavenumber — far more expensive than
FFTs.

\subsection{The Pseudo-Spectral Solution}

The \textbf{pseudo-spectral trick}: compute products in physical space
(cheap, pointwise), and everything else in Fourier space:
\begin{enumerate}
  \item Transform the fields you need from Fourier space to physical space (IFFT).
  \item Compute the product pointwise (simple array multiplication).
  \item Transform the result back to Fourier space (FFT).
\end{enumerate}
Each FFT/IFFT costs $O(N^2\log N)$ in 2D.  The total cost is much lower than
convolution.

\subsection{The 5-Step RHS Pipeline}

Every time the ODE solver requests $d\hat\theta/dt$, the following five steps
run.

\paragraph{Step 1 — Compute $g = \theta\,\nabla\lapinv\theta$ (pseudo-spectral).}
\begin{enumerate}[label=\alph*.]
  \item Apply $\lapinv$ in Fourier space:
    $\hat\psi = \LAPINV\cdot\hat\theta$.
  \item Compute $\psi = \mathrm{real}(\mathrm{IFFT2}(\hat\psi))$ (physical
    space stream-function-like field).
  \item Compute derivatives:
    $\partial_x\psi = \mathrm{real}(\mathrm{IFFT2}(\DELX\cdot\hat\psi))$ and
    $\partial_y\psi$ similarly.
  \item Compute the nonlinear product in physical space:
    $g_x = \theta\cdot\partial_x\psi$, $g_y = \theta\cdot\partial_y\psi$.
  \item Transform back: $\hat g_x = \mathrm{FFT2}(g_x)$,
    $\hat g_y = \mathrm{FFT2}(g_y)$.
\end{enumerate}

\paragraph{Step 2 — Apply Leray projection (pure Fourier-space multiplication).}
\[
  (\hat g_x,\hat g_y)\;\longrightarrow\;(\widehat{Pg_x},\widehat{Pg_y}).
\]

\paragraph{Step 3 — Apply $-\lapinv$ to get velocity $v = -\lapinv Pg$.}
\[
  \hat v_x = -\LAPINV\cdot\widehat{Pg_x}, \qquad
  \hat v_y = -\LAPINV\cdot\widehat{Pg_y}.
\]

\paragraph{Step 4 — Normalise to enforce the enstrophy constraint $\|\nabla u\|_{L^2} = F$ (via Parseval).}
\[
  \|\nabla v\|_{L^2}
  = \frac{1}{N^2}\sqrt{\|\DELX\hat v_x\|^2+\|\DELY\hat v_x\|^2
                     +\|\DELX\hat v_y\|^2+\|\DELY\hat v_y\|^2},
  \qquad
  \hat u = F\,\frac{\hat v}{\|\nabla v\|_{L^2}}.
\]

\paragraph{Step 5 — Advect to compute $d\hat\theta/dt = -\mathcal{F}[u\cdot\nabla\theta]$.}
\[
  \frac{d\hat\theta}{dt}
  = -\mathrm{FFT2}\!\left[u_x\cdot\mathrm{IFFT2}(\DELX\cdot\hat\theta)
    + u_y\cdot\mathrm{IFFT2}(\DELY\cdot\hat\theta)\right].
\]

% ─────────────────────────────────────────────────────────────────────────────
\section{Part 11: Time Integration with RK45}

\subsection{The ODE}

After spectral discretisation, the PDE becomes a large system of ODEs:
\begin{equation}
  \frac{d\hat\theta}{dt} = F(\hat\theta,t),
  \qquad \hat\theta(0) = \widehat{\theta_0},
\end{equation}
where $F$ is the RHS from the five steps above and $\hat\theta$ has $N^2$
complex components.

\subsection{What Is RK45?}

The \textbf{Runge--Kutta 4/5} (Dormand--Prince) method is a classical
adaptive ODE solver.  Here is the idea:
\begin{itemize}
  \item At each time step $t_n$, evaluate $F$ at several cleverly chosen
    intermediate points between $t_n$ and $t_n + h$.
  \item Combine these evaluations to get a 4th-order \emph{and} a 5th-order
    estimate of $\hat\theta(t_n+h)$.
  \item Use the difference between these two estimates to measure the local
    error.
  \item If the error is too large, halve the step size $h$ and retry; if it
    is very small, increase $h$ to speed up.
\end{itemize}
This is called an \textbf{adaptive} time integrator — it automatically picks a
step size that balances accuracy and speed.

\paragraph{Why RK45 and not simpler methods?}
Euler's method (first order) is too inaccurate.  Fourth-order Runge--Kutta
with a fixed step is accurate but wastes evaluations when the solution is
smooth.  RK45 is the standard workhorse for stiff-to-moderate ODEs and matches
MATLAB's \code{ode45}.

\begin{lstlisting}
sol = scipy.integrate.solve_ivp(
    rhs,
    t_span=(0, T),
    y0=theta_hat_flat_real,
    method='RK45',
    rtol=1e-6,
    atol=1e-8
)
\end{lstlisting}

\subsection{Real/Imaginary Splitting}

\code{solve\_ivp} requires a real-valued state vector.  The complex array
$\hat\theta\in\mathbb{C}^{N^2}$ is stored as:
\[
  y = \bigl[\mathrm{Re}(\hat\theta_{\mathrm{flat}}),\;
            \mathrm{Im}(\hat\theta_{\mathrm{flat}})\bigr]
    \in \mathbb{R}^{2N^2}.
\]
At each RHS call, the code reconstructs $\hat\theta$ from $y$, performs the
five-step computation, and returns $dy/dt$ in real form.

% ─────────────────────────────────────────────────────────────────────────────
\section{Part 12: The Resolution Check — When to Stop}

\subsection{Why the Simulation Must Stop}

As the velocity stirs $\theta$ into finer and finer filaments, the scalar
develops features at increasingly high wavenumbers.  Eventually these features
become smaller than the grid spacing $dx = 1/N$ — the grid can no longer
represent them.  Beyond this point, the simulation is no longer physically
meaningful.

\subsection{What Is Conserved?}

Under exact incompressible advection, \textbf{all $L^p$ norms are conserved}:
\[
  \|\theta(t)\|_{L^p} = \|\theta_0\|_{L^p}
  \quad \text{for all } t\geq 0,\ \text{all } p\geq 1.
\]
This is a mathematical fact following from the change-of-variables formula for
divergence-free flows.  Numerically, these norms should remain constant as long
as the grid can resolve the solution.

\subsection{The Event Function}

The code monitors:
\begin{equation}
  e(t) = \max\!\left(
    \bigl|\|\theta\|_{L^2}-1\bigr|,\;
    \left|\frac{\|\theta\|_{L^4}}{\|\theta_0\|_{L^4}}-1\right|,\;
    \left|\frac{\|\theta\|_{L^8}}{\|\theta_0\|_{L^8}}-1\right|
  \right) - \varepsilon,
  \qquad \varepsilon = 10^{-3}.
\end{equation}
When $e(t) > 0$ (any norm drifts by more than $0.1\%$), the solver stops.

\paragraph{Why also use $L^4$ and $L^8$?}
The $L^2$ norm is quadratic and relatively forgiving.  Higher $L^p$ norms are
more sensitive to large values and sharp gradients — exactly the kind of
structures that appear when the grid is under-resolving.  Using three norms
gives an earlier and more reliable warning.

% ─────────────────────────────────────────────────────────────────────────────
\section{Recommended Books and Resources}

\subsection*{Introductory — Start Here}

\begin{description}[leftmargin=0pt]
  \item[\normalfont\textbf{``A Student's Guide to Fourier Transforms''}]
    J.\,F.\ James, Cambridge University Press, 3rd~ed., 2011.\\
    The most accessible introduction.  Written specifically for science
    students without a heavy maths background.  Short chapters, physical
    intuition, minimal prerequisites.

  \item[\normalfont\textbf{``The Fourier Transform and Its Applications''}]
    Ronald Bracewell, McGraw-Hill, 3rd~ed., 2000.\\
    A classic with excellent hand-drawn visualisations.  Bracewell's ``picket
    fence'' picture of the DFT and his sampling-theorem explanations are
    famous for being unusually clear.

  \item[\normalfont\textbf{``Fourier Analysis: An Introduction''}]
    Elias Stein \& Rami Shakarchi, Princeton Lectures in Analysis, Vol.~1,
    2003.\\
    More rigorous, but very clearly written.  Good for understanding
    \emph{why} the theorems hold, not just how to use them.
\end{description}

\subsection*{Spectral / Numerical Methods — Directly Relevant}

\begin{description}[leftmargin=0pt]
  \item[\normalfont\textbf{``Spectral Methods in MATLAB''}]
    Lloyd N.\ Trefethen, SIAM, 2000.  \emph{Free PDF online.}\\
    The perfect companion to this project.  40~short programs (easily
    translated from MATLAB to Python) explaining spectral methods for PDEs.
    Very readable; starts from scratch.  This book will make Section~3 of the
    report completely transparent.

  \item[\normalfont\textbf{``Numerical Methods for Engineers''}]
    Chapra \& Canale, McGraw-Hill, 7th~ed., 2015.\\
    Covers FFT, ODE solvers (including RK45), and numerical differentiation at
    an introductory engineering level.  Very practical and hands-on.

  \item[\normalfont\textbf{``A Practical Guide to Pseudospectral Methods''}]
    Bengt Fornberg, Cambridge University Press, 1996.\\
    Specifically about pseudo-spectral methods — exactly the approach used in
    the mixing simulation.  More technical than Trefethen's book, but the most
    thorough treatment of the subject.

  \item[\normalfont\textbf{``Numerical Methods for Fluid Dynamics''}]
    Dale Durran, Springer, 2nd~ed., 2010.\\
    Covers spectral methods, pseudo-spectral methods, and time integration at
    a graduate level.  Bridges the gap between numerical methods and physical
    simulation.
\end{description}

\subsection*{PDEs and the Physical Problem}

\begin{description}[leftmargin=0pt]
  \item[\normalfont\textbf{``Introduction to Partial Differential Equations''}]
    Walter Strauss, Wiley, 2nd~ed., 2008.\\
    Good introductory PDE text covering Fourier series and transforms from a
    mathematical perspective.  Accessible to anyone with calculus.

  \item[\normalfont\textbf{``An Introduction to Fluid Dynamics''}]
    G.\,K.\ Batchelor, Cambridge University Press, 2000.\\
    The definitive fluid dynamics reference.  Chapters~1--3 are accessible and
    explain divergence-free velocity fields, the incompressibility condition,
    and passive scalar transport.
\end{description}

\subsection*{Online Resources}

\begin{center}
\begin{tabular}{ll}
  \toprule
  Resource & What It Covers \\
  \midrule
  3Blue1Brown ``But What Is the Fourier Transform?'' (YouTube)
    & Best 20-min visual intro to Fourier transforms \\
  3Blue1Brown ``But What Is a Fourier Series?'' (YouTube)
    & Visual intuition for decomposing functions into waves \\
  NumPy FFT documentation
    & Exact conventions used in \code{np.fft.fft2} \\
  Trefethen's ATAP (free PDF)
    & Rigorous approximation theory, spectral methods \\
  SciPy \code{solve\_ivp} documentation
    & Full details of the RK45 implementation \\
  \bottomrule
\end{tabular}
\end{center}

% ─────────────────────────────────────────────────────────────────────────────
\section{Summary Table}

\begin{center}
\renewcommand{\arraystretch}{1.25}
\begin{tabular}{lll}
  \toprule
  Concept & What It Is & Code Equivalent \\
  \midrule
  DFT
    & Decompose discrete function into frequencies
    & \code{np.fft.fft2} \\
  FFT
    & Fast $O(N\log N)$ algorithm for DFT
    & \code{np.fft.fft2} \\
  Wavenumber $k$
    & Frequency index; oscillations per unit length
    & \code{k\_eff} \\
  Centred wavenumbers
    & Remap DFT indices to $\{-N/2,\ldots,N/2\}$
    & \code{k - N*(k > N//2)} \\
  Nyquist frequency
    & Maximum representable frequency $= N/2$
    & Index \code{N//2} \\
  Spectral derivative
    & Multiply $\hat f$ by $2\pi ik$
    & \code{DEL\_X}, \code{DEL\_Y} \\
  Inverse Laplacian
    & Multiply $\hat f$ by $-1/(4\pi^2\|k\|^2)$
    & \code{LAP\_INV} \\
  Parseval identity
    & $\|f\|_{L^2} = \|\hat f\|_{\ell^2}/N^2$ (2D)
    & Norm computations \\
  $H^{-1}$ mix norm
    & Weighted norm that decreases with mixing
    & \code{LAMBDA\_INV} \\
  Leray projection
    & Extracts divergence-free part of vector field
    & Step~2 of RHS \\
  Pseudo-spectral
    & Products in physical, derivatives in Fourier
    & RHS evaluation \\
  RK45
    & Adaptive-step 4/5 Runge--Kutta ODE solver
    & \code{solve\_ivp} \\
  Resolution check
    & Stop when $L^p$ norms drift beyond tolerance
    & \code{make\_res\_check()} \\
  \bottomrule
\end{tabular}
\end{center}

\bigskip
\noindent\textit{This guide accompanies the notebook
\code{00\_dft\_for\_beginners.ipynb} in the \code{python\_code/} folder,
which provides interactive Python demonstrations of every concept described
here.}

\end{document}
